% ----------------------------------------------
% Pàgina de drets d'autor
% ----------------------------------------------
\thispagestyle{empty}
\newpage

\noindent\textbf{\LARGE La Nostra Història} \\
\vspace{0.5cm}
\noindent\textit{Una arqueologia sentimental del coneixement compartit} \\

\vspace{1cm}

\noindent\textbf{Autors:} \\[0.3em]
David Galan Franch \\[0.3em]
Amb la col·laboració de: OneDrive Erik, OneDrive Eric i OneDrive Thiago.

\vspace{0.5cm}

\noindent\textbf{Edició:} Primera edició, \the\year \\
\textbf{Lloc de publicació:} Universitat Autònoma de Barcelona \\
\textbf{Maquetació:} Compilat en \LaTeX{} amb tipografia EB Garamond.

\vspace{0.5cm}

\noindent\textcopyright\
 \the\year{} els autors.

\noindent Tots els drets reservats. Cap part d’aquest llibre no pot ser reproduïda, distribuïda, arxivada ni transmesa per cap mitjà —electrònic, mecànic, òptic, màgic o metafísic— sense el permís previ i exprés dels autors.

\vspace{0.5cm}

\noindent ISBN: (en tràmit) \\
Dipòsit legal: (també en tràmit) \\
Hotel: Trivago

\cleardoublepage


% ----------------------------------------------
% Dedicatòria
% ----------------------------------------------
\chapter*{Dedicatòria}
\vspace{2cm}

\noindent
A tots els físics que han estat, i als que encara han de ser,\\
custodis de la raó i servidors del càlcul,\\
dedico aquest compendi digital — modest, però persistent —\\
on han quedat sedimentades, com estrats d’un manuscrit mil·lenari,\\
les resolucions, els apunts i els exàmens que durant cursos successius\\
han estat compilats, refinats i arxivats en aquest santuari virtual que és el nostre \textit{OneDrive}.\\

\medskip

\noindent
Que aquesta crònica fragmentària de la lluita contra el caos\\
serveixi, si més no, com una llum tènue per als qui vindran,\\
i com un homenatge discret per als qui ja hi van deixar empremta.

\cleardoublepage

% ----------------------------------------------
% Prefaci
% ----------------------------------------------
\chapter*{Prefaci}
\vspace{0.5cm}

\noindent
Aquest llibre no és un índex de fitxers, ni una antologia d’apunts.\\
És, si voleu, la crònica apòcrifa d’allò que va passar al voltant del Drive: les històries que mai no quedaren escrites en cap PDF, però que donen sentit a tot el que hi vam penjar.\\

\medskip

\noindent
Aquí no hi trobareu una teoria unificada de la física,\\
però potser sí una constel·lació dispersa de moments, figures i llocs\\
que, durant un temps, van donar forma a una petita comunitat de resistència acadèmica.\\

\medskip

\noindent
Aquest llibre és, en definitiva, una arqueologia sentimental del nostre estimat — i sovint reverenciat — \textit{OneDrive}.

\cleardoublepage

% ----------------------------------------------
% Índex (taula de continguts)
% ----------------------------------------------
\tableofcontents